\documentclass[11pt,a4paper]{article}
\usepackage[T1]{fontenc}
\usepackage{amsmath,amssymb,amsthm,geometry,booktabs,hyperref,longtable}
\usepackage[htt]{hyphenat}
\emergencystretch=3em
\sloppy
\geometry{margin=2.5cm}
\newtheorem{finding}{Finding}
\theoremstyle{remark}
\newtheorem*{remark}{Remark}
\newcommand{\LG}{\mathrm{LG}}
\newcommand{\OG}{\mathrm{OG}}
\newcommand{\Dg}{D}
\title{\textbf{Peer review: Theorem I and Theorem I\textsubscript{OG}}\\[4pt]
\large the integer zero window, exactly, on $\LG(m,2m)$ and $\OG(m,2m{+}1)$\\[2pt]
\normalsize \emph{endorsed; with a shorter proof of its load-bearing lemma,
and a correction to my own closed form}}
\author{Clio (clio-vega)}
\date{2026-10-24}

\makeatletter
%% MIRROR of proofs/protocol-header.tex (clio-vega/proofs), copied 2026-10-16.
%% reviews/ is a SEPARATE REPO (clio-vega/rick-review) and must build standalone for
%% Rick and Robin, so a symlink or TEXINPUTS dependency would break for them. Hence a
%% copy -- but a VERIFIABLE one, because a silent fork is the same fault in new clothes:
%%   source md5 at copy time: f828569e31971b23bc5726db5d23ecef
%%   drift check:  md5sum ../proofs/protocol-header.tex   # compare with the line above
%% If they differ, re-copy from proofs/ -- that file is the source of truth.
%%
%% protocol-header.tex -- PROTOCOL §2.3 page-1 block, as a TEMPLATE not an instance.
%%
%% PROTOCOL §2.3: "Every PDF must state, on its first page: author, date,
%% recipient, a title, and the commit hash of the work-in-progress push it
%% corresponds to.  A PDF with no commit hash cannot be checked against the
%% record and should be treated as a draft opinion, not a result."
%%
%% Written 2026-10-16 because the fault had recurred twice and I had twice
%% fixed only the instance:
%%   - 2026-10-14: the Q473 PDF went to Rick with no recipient and no hash.
%%   - 2026-10-14: the Q460 PDF went to Rick saying "Recipient: Robin Langer".
%%     He noticed, and took the commit hash as authority instead.
%% Both times I repaired the one file and wrote "fix the template, not just the
%% instance" in my notes.  This file is that fix.  \input it; do not retype it.
%%
%% USAGE, immediately after \maketitle:
%%   %% MIRROR of proofs/protocol-header.tex (clio-vega/proofs), copied 2026-10-16.
%% reviews/ is a SEPARATE REPO (clio-vega/rick-review) and must build standalone for
%% Rick and Robin, so a symlink or TEXINPUTS dependency would break for them. Hence a
%% copy -- but a VERIFIABLE one, because a silent fork is the same fault in new clothes:
%%   source md5 at copy time: f828569e31971b23bc5726db5d23ecef
%%   drift check:  md5sum ../proofs/protocol-header.tex   # compare with the line above
%% If they differ, re-copy from proofs/ -- that file is the source of truth.
%%
%% protocol-header.tex -- PROTOCOL §2.3 page-1 block, as a TEMPLATE not an instance.
%%
%% PROTOCOL §2.3: "Every PDF must state, on its first page: author, date,
%% recipient, a title, and the commit hash of the work-in-progress push it
%% corresponds to.  A PDF with no commit hash cannot be checked against the
%% record and should be treated as a draft opinion, not a result."
%%
%% Written 2026-10-16 because the fault had recurred twice and I had twice
%% fixed only the instance:
%%   - 2026-10-14: the Q473 PDF went to Rick with no recipient and no hash.
%%   - 2026-10-14: the Q460 PDF went to Rick saying "Recipient: Robin Langer".
%%     He noticed, and took the commit hash as authority instead.
%% Both times I repaired the one file and wrote "fix the template, not just the
%% instance" in my notes.  This file is that fix.  \input it; do not retype it.
%%
%% USAGE, immediately after \maketitle:
%%   %% MIRROR of proofs/protocol-header.tex (clio-vega/proofs), copied 2026-10-16.
%% reviews/ is a SEPARATE REPO (clio-vega/rick-review) and must build standalone for
%% Rick and Robin, so a symlink or TEXINPUTS dependency would break for them. Hence a
%% copy -- but a VERIFIABLE one, because a silent fork is the same fault in new clothes:
%%   source md5 at copy time: f828569e31971b23bc5726db5d23ecef
%%   drift check:  md5sum ../proofs/protocol-header.tex   # compare with the line above
%% If they differ, re-copy from proofs/ -- that file is the source of truth.
%%
%% protocol-header.tex -- PROTOCOL §2.3 page-1 block, as a TEMPLATE not an instance.
%%
%% PROTOCOL §2.3: "Every PDF must state, on its first page: author, date,
%% recipient, a title, and the commit hash of the work-in-progress push it
%% corresponds to.  A PDF with no commit hash cannot be checked against the
%% record and should be treated as a draft opinion, not a result."
%%
%% Written 2026-10-16 because the fault had recurred twice and I had twice
%% fixed only the instance:
%%   - 2026-10-14: the Q473 PDF went to Rick with no recipient and no hash.
%%   - 2026-10-14: the Q460 PDF went to Rick saying "Recipient: Robin Langer".
%%     He noticed, and took the commit hash as authority instead.
%% Both times I repaired the one file and wrote "fix the template, not just the
%% instance" in my notes.  This file is that fix.  \input it; do not retype it.
%%
%% USAGE, immediately after \maketitle:
%%   \input{protocol-header}
%%   \protocolheader{Rick (grandpa-rick)}{clio-vega/proofs}{60e4ce3}
%%
%% The recipient argument is MANDATORY and has no default on purpose: a default
%% is what produced "Recipient: Robin Langer" on a PDF addressed to Rick.

%% \maketitle CLEARS \@date (and \@author, \@title) in the article class, so a
%% macro called AFTER \maketitle that reads \@date prints an EMPTY row.  Caught
%% 2026-10-16 only because I ran pdftotext on page 1 instead of trusting rc=0:
%% the build was clean and the Date field was blank.  Hence: snapshot the date
%% HERE, at \input time in the preamble, where \@date is still live.
\let\protocol@date\@date

\newcommand{\protocolheader}[3]{%
  \begin{center}
  \fbox{\begin{minipage}{0.92\textwidth}
  \small
  \textbf{PROTOCOL §2.3 record.}\\[1mm]
  \begin{tabular}{@{}ll@{}}
  Author     & Clio \\
  Date       & \protocol@date \\
  Recipient  & #1 \\
  Repository & \texttt{#2} \\
  Commit     & \texttt{#3} \\
  \end{tabular}\\[1mm]
  This PDF corresponds to the pushed commit above.  Where this text and the
  commit disagree, \textbf{the commit is authoritative}.
  \end{minipage}}
  \end{center}
  \vspace{2mm}%
}

%%   \protocolheader{Rick (grandpa-rick)}{clio-vega/proofs}{60e4ce3}
%%
%% The recipient argument is MANDATORY and has no default on purpose: a default
%% is what produced "Recipient: Robin Langer" on a PDF addressed to Rick.

%% \maketitle CLEARS \@date (and \@author, \@title) in the article class, so a
%% macro called AFTER \maketitle that reads \@date prints an EMPTY row.  Caught
%% 2026-10-16 only because I ran pdftotext on page 1 instead of trusting rc=0:
%% the build was clean and the Date field was blank.  Hence: snapshot the date
%% HERE, at \input time in the preamble, where \@date is still live.
\let\protocol@date\@date

\newcommand{\protocolheader}[3]{%
  \begin{center}
  \fbox{\begin{minipage}{0.92\textwidth}
  \small
  \textbf{PROTOCOL §2.3 record.}\\[1mm]
  \begin{tabular}{@{}ll@{}}
  Author     & Clio \\
  Date       & \protocol@date \\
  Recipient  & #1 \\
  Repository & \texttt{#2} \\
  Commit     & \texttt{#3} \\
  \end{tabular}\\[1mm]
  This PDF corresponds to the pushed commit above.  Where this text and the
  commit disagree, \textbf{the commit is authoritative}.
  \end{minipage}}
  \end{center}
  \vspace{2mm}%
}

%%   \protocolheader{Rick (grandpa-rick)}{clio-vega/proofs}{60e4ce3}
%%
%% The recipient argument is MANDATORY and has no default on purpose: a default
%% is what produced "Recipient: Robin Langer" on a PDF addressed to Rick.

%% \maketitle CLEARS \@date (and \@author, \@title) in the article class, so a
%% macro called AFTER \maketitle that reads \@date prints an EMPTY row.  Caught
%% 2026-10-16 only because I ran pdftotext on page 1 instead of trusting rc=0:
%% the build was clean and the Date field was blank.  Hence: snapshot the date
%% HERE, at \input time in the preamble, where \@date is still live.
\let\protocol@date\@date

\newcommand{\protocolheader}[3]{%
  \begin{center}
  \fbox{\begin{minipage}{0.92\textwidth}
  \small
  \textbf{PROTOCOL §2.3 record.}\\[1mm]
  \begin{tabular}{@{}ll@{}}
  Author     & Clio \\
  Date       & \protocol@date \\
  Recipient  & #1 \\
  Repository & \texttt{#2} \\
  Commit     & \texttt{#3} \\
  \end{tabular}\\[1mm]
  This PDF corresponds to the pushed commit above.  Where this text and the
  commit disagree, \textbf{the commit is authoritative}.
  \end{minipage}}
  \end{center}
  \vspace{2mm}%
}

\makeatother

\begin{document}
\maketitle
\protocolheader{Rick (grandparick20@gmail.com), cc Robin}{clio-vega/rick-review}{297d85a1f93d4a31fa9fa22dcd47c1e6a5573275}

\begin{center}
\begin{tabular}{p{3.0cm}p{11.2cm}}
\toprule
\textbf{Target} & \texttt{proofs/2026-10-24-integer-window-and-OG-twin.md} (308 lines) \\
\textbf{Reviewed at} & \texttt{rick-research} \texttt{7ff2361d182d258dbf6c829c4aa97bcfd9fffc8a} \\
\textbf{Arming stamp} & \texttt{7bd9ae3c}; \texttt{origin/main} had moved to \texttt{b44a146}
  (11 commits) \\
\textbf{Covering mail} & UID 847; context UIDs 848, 850, 851 \\
\textbf{Verification code} & \texttt{clio-vega/rick-review} \texttt{297d85a},
  dir \texttt{code-1024c2-thmI/} \\
\bottomrule
\end{tabular}
\end{center}

\section*{0. First instrument: the commit range}

Eleven commits landed between the arming stamp and this session. The target file did
\emph{not} move, so I reviewed the version the brief named. Three things in the range matter:
PROVE/Dream 254 upgraded many registry nodes \texttt{computed}$\to$\texttt{proved} (including
\textbf{Theorem B}, which I had left unendorsed, \S6); \texttt{0e4001a} gives a BR-free second
proof of SR\textsubscript{LG}; and the new file
\texttt{proofs/2026-10-25-crownF-strict-reciprocity.md} contains \textbf{Theorem Q30}, which
\emph{retires the reader-left dependency Theorem I still declares for itself}. That is
Finding~2, and it moves the grade upward.

Inbox re-checked immediately before sending, not at arming: UIDs 849, 850, 851 had landed past
the watermark and are answered in \S8.

\section*{1. The claim, in one sentence}

With $s_j(\kappa) := (-1)^{|\kappa|}H_\kappa(-j)$, one monotone induction on $|\kappa|$ ---
driven by the K-theoretic recursion
\[
  s_{j+1}(\kappa) \;=\; s_j(\kappa) \;+\; \sum_\rho |c_\rho|\, s_j(\rho)
  \tag{$\star$}
\]
whose increments are all $\ge 0$ \emph{because the anchor $s_1=0$ leaves the alternating signs
nothing to fight} --- yields the exact integer zero set $\{-1,\dots,-W\}$ with
$W=\lceil r(\kappa)/2\rceil$ on $\LG$ and $W=r(\kappa)$ on $\OG$, plus the sign
$(-1)^{|\kappa|}$ beyond the window.

The mechanism is the best idea in the file, and \S11 names it correctly: when an
alternating-sign recursion refuses to be coherent, look for the specialisation where the base
case vanishes.

\section*{2. Verdict}

\begin{center}
\small
\begin{tabular}{p{3.0cm}p{3.5cm}p{6.4cm}}
\toprule
\textbf{Item} & \textbf{Grade I endorse} & \textbf{Basis} \\
\midrule
Lemma 0 (chain formula) & proved & 510 $\kappa$, $m\le8$, three definitions agree, 0 mismatches \\
Lemma 1 (anchor $H_\kappa(-1)=0$) & proved given (G1) & read, incl.\ Lemma P which licenses it \\
Lemma 2 (Pieri support, sign) & proved given (BR) & every quote verbatim, every locator correct \\
Lemma Ch$'$ (chain drop) & proved & citation checked; 0 violations, drops exactly $2/1$, $m\le7$ \\
Lemma A (i),(ii),(iv),(v) & proved (one scope slip, F1) & 502 $\kappa$, $m\le8$ \\
\textbf{Lemma A(iii)} \emph{(sole nonvanishing witness)} & proved given (BR) & \textbf{502/502}, $m\le8$; two-line proof in \S7 \\
\textbf{Theorem I ($\LG$)} & \textbf{proved given (G0),(G1),(BR)} --- ENDORSED & independent engine $m\le6$, 0 failures, 8 control arms firing \\
\textbf{Theorem I\textsubscript{OG}} & \textbf{proved given (G0),(G1),(BR)} --- ENDORSED & independent \textbf{type-B} engine, $m\le6$, 0 failures \\
Cor NZ\textsubscript{R} (\S12) & proved given (G0),(G1),(BR) & conditional on Prop.\ R, which I did not re-read \\
R1/R2, Lemma E (\S12) & computed; gap self-declared, correctly scoped & not independently verified \\
\bottomrule
\end{tabular}
\end{center}

\medskip
\noindent\textbf{The (BR) debt is lighter than the file says.} \S10 gap~2 declares a dependency
on BR Cor~5.6, whose proof BR leave to the reader. His own Theorem Q30 derives that statement
from BR Thm~5.5, which BR prove in full. I verified Theorem Q30 independently,
\textbf{3025/3025}. So the honest grade is \texttt{proved given (G0),(G1),(BR)} with
\textbf{(BR) now meaning only results BR prove}.

\section*{3. The pre-registered falsifier: P1 CONFIRMED}

P1 asked whether the \emph{strict} step $s_W<s_{W+1}$ rests on a nonzero witness or on
``$\mathrm{supp}=\mathrm{rims}$''. The latter would be unsound: my 10-23 census found
$1252/3025$ rim shapes carrying coefficient $0$, re-measured here at \textbf{41.4\%}.

\textbf{It rests on a nonzero witness.} Lemma A(iii) \emph{exhibits} a KLG-tableau of shape
$R(T_{t^*})$. Nothing appeals to $\mathrm{supp}=\mathrm{rims}$. Better: Lemma 2 writes
$c_\rho=(-1)^{|\theta|-1}N(\theta)$ with $N(\theta)\ge0$, so a vanishing coefficient on a rim is
\emph{built into the statement}, and \S6's control records 42 of 67 $r$-odd 2-strips vanishing.
He knew the correction and used the right version in the right place. I verified the witness
rather than trusting the construction: $N(T_{t^*})\ge1$ on \textbf{502/502} nonempty $\kappa$ at
$m\le8$. The vanishing half needs only the containment, so the zero rims are harmless there.
Both halves are correctly scoped.

\section*{4. Verified by reading}

\textbf{Lemma P is the premise the anchor really rests on}, because Lemma 1 evaluates the
Demazure expression at the \emph{non-dominant} weight $j=-1$. I worked through it: the operator
identity $D_if=(f-s_if)/(1-e^{-\alpha_i})+s_if$ is equivalent to the standard form; $f-s_if$
vanishes on $\alpha_i=0$ since $s_i$ fixes that hyperplane pointwise, and
$1-e^{-\alpha_i}=\alpha_i\cdot(\text{unit})$, so $D_i(\mathfrak{m}^k)\subseteq\mathfrak{m}^{k-1}$;
hence $D_w(\mathfrak{m}^{L+1})\subseteq\mathfrak{m}$, the Taylor tail drops, and $P_w(\lambda)$
is polynomial of degree $\le\ell(w)$. \textbf{Sound --- the anchor is safe.} This was the premise
I most expected to break.

\textbf{Lemma Ch$'$'s cross-file citation is accurate.} \S5 claims the Z file's Lemma Ch proof
``never uses row 0''. I read it: the three-box case analysis uses only the rim condition and the
$\le2$ bound; only the \emph{consequence} paragraph used $\Dg(\kappa)_+$, and Lemma Ch$'$
correctly re-derives it on the full diagram via Lemma 0. A verified locator, not a vibe.

\textbf{(G0) is invoked, not merely listed} (my 10-24 F1 asked for this): it is consumed at the
Recursion step, $b_\kappa(j)=\chi([\mathcal O_{X_\kappa}]\cdot(1-[L^{-1}])\cdot[L^j])$. I also
re-derived $(\star)$ by hand; the sign bookkeeping
($c_\rho=(-1)^{|\theta|-1}|c_\rho|$, $|\kappa|=|\rho|+|\theta|$, $(-1)^{2|\theta|-1}=-1$) is
correct, including the cancellation that leaves $+\sum|c_\rho|s_j(\rho)$.

\textbf{Every BR citation checks out, first-hand}, against \texttt{1005.2605}~v1 source on my
disk, with the theorem counter reconstructed (\verb|\newtheorem{thm}{Theorem}[section]|, shared
by lemma/prop/cor):

\begin{center}
\small
\begin{tabular}{p{5.0cm}p{4.0cm}p{3.6cm}}
\toprule
\textbf{His label} & \textbf{Actual number} & \textbf{Check} \\
\midrule
``unlabelled Cor after eq.\ (hsumB)'' & \textbf{Cor 4.4}, line 956 (eq.\ line 948) & verbatim \\
Lemma \texttt{vanishB} & \textbf{Lemma 4.5}, line 961 & (a),(b),(c) verbatim \\
Cor \texttt{pieriCtab} & \textbf{Cor 5.6}, line 1567 & verbatim \\
rook strip $=\le1$ per row/col, ``\S3'' & line 537, inside \S3 & locator correct \\
``OG, \S4'' / ``LG, \S5'' & \S4 max.\ orthogonal / \S5 Lagrangian & correct \\
KLG conditions (i)--(v), $\{1'{<}1{<}2'{<}2{<}\cdots\}$ & lines 1553--1562 & all five verbatim \\
``left to the reader'' for Cor 5.6 & line 1564 & verbatim \\
\bottomrule
\end{tabular}
\end{center}

\noindent No phantom determinants. One subtlety he handled correctly: BR's KLG definition opens
``Let $\theta$ be a \emph{rim}'', so the count is defined only on rims --- exactly the support
condition Lemma 2 imposes.

\section*{5. Independent computation}

The brief required a positive control \emph{including the OG side}; my 10-24 census was LG-only.
I built a \textbf{type-B $(\mathrm{SO}(2m{+}1))$ Demazure engine} (weights stored as $2\lambda$
so the spinor weight stays integral) and validated it before use: $H^{\OG}_\kappa(0)=1$ for all
$\kappa$, and $H^{\OG}_{(2,1)}(j)=\binom{j+3}{3}$, the order polynomial of the 3-chain
$\Dg(2,1)$.

\begin{center}
\small
\begin{tabular}{p{4.1cm}p{4.4cm}p{4.6cm}}
\toprule
\textbf{Check} & \textbf{Domain} & \textbf{Result} \\
\midrule
Lemma 0, three definitions & 510 $\kappa$, $m\le8$ & 0 mismatches \\
Lemma A (i),(ii),(iv),(v) & 502 $\kappa$, $m\le8$ & 0 failures but F1 \\
\textbf{Lemma A(iii)} $N(T_{t^*})\ge1$ & 502 $\kappa$, $m\le8$ & \textbf{502/502} \\
Lemma Ch$'$ drop $\le2$ / $\le1$ & true support, $m\le7$ (LG 1168, OG \textbf{449}) & 0 violations; max exactly $2/1$ \\
\textbf{Theorem I} window & $m\le6$, 126 $\kappa$ & vanishing 359/359, positive 630/630, monotone 863/863 \\
\textbf{Theorem I\textsubscript{OG}} window & $m\le6$, 126 $\kappa$ & vanishing 645/645, positive 630/630, monotone 1149/1149 \\
Theorem Q30 closed form & 3025 rim pairs, $m\le7$ & \textbf{3025/3025} \\
Theorem B conclusion (\S6) & 150 $(\kappa,a)$, $m\le4$ & \textbf{150/150} \\
\bottomrule
\end{tabular}
\end{center}

My OG support count \textbf{449 at $m=7$ matches his 449 exactly}. The LG figures separate
usefully from his: \textbf{1337} admissible pairs but only \textbf{1168} with nonzero
coefficient --- the 169-pair difference \emph{is} the support correction, measured.

\textbf{Which $\kappa$}: all 126 strict $\kappa$ with parts in $[m]$ for $m\le6$, i.e.\
\emph{every} $\kappa$, not a sample --- including $r(\kappa)$ up to 11, both parities of $r$, and
the staircase $(6,5,4,3,2,1)$. A census domain can be the already-proved class; this one is not.

\paragraph{Negative controls --- every arm fired.}
\begin{center}
\small
\begin{tabular}{p{5.0cm}p{2.6cm}p{5.2cm}}
\toprule
\textbf{Control} & \textbf{Expected} & \textbf{Result} \\
\midrule
window shifted $\pm1$ (8 arms, both types) & must fire & \textbf{8/8 fire} \\
use $t=2$ when $r$ odd (wrong parity) & must fail & $N(T_2)=0$ on \textbf{163/293} \\
$\mathrm{supp}=\mathrm{rims}$ & must be false & \textbf{41.4\%} of rims carry 0 \\
Theorem B without $a\ge m$ & must fail & \textbf{26/68} violate $\ge1$ \\
closed form on $\theta$ instead of $R(\theta)$ & must fire & \textbf{528/3025} wrong \\
\midrule
\emph{positive}: \texttt{rho\_of\_ideal}$(\Dg(\kappa))=\kappa$ & must pass & 510/510 \\
\emph{positive}: closed form vs brute force & must pass & 2033 shapes, 0 mismatches \\
\bottomrule
\end{tabular}
\end{center}

\begin{remark}[how I caught my own bug]
My first Lemma A run reported \textbf{2963 failures}, and I nearly had a finding. But Lemma 0
and Lemma Ch$'$ were simultaneously clean, and clause (i) ``the complement is an order ideal''
--- near-tautological --- failed on 453 of 502. A theorem does not fail that way; an instrument
does. \texttt{rho\_of\_ideal} was adding a spurious $+i$ to each part. After repair: 8 failures,
all explained (F1). \emph{The asymmetry between a failing clause and its clean neighbours was
the diagnostic, not the count.}
\end{remark}

\paragraph{A by-product: his Q33, answered by his own Lemma P.}
\S8 records \texttt{check.py 5} and \texttt{mono.py 5} \textbf{killed} (50 min, exact-Fraction
interpolation at degree 15), so $m\ge5$ is unchecked in the file. But Lemma P already licenses
direct evaluation at negative weights: for integral $\lambda$ the germ evaluation at $h=0$
\emph{is} the Laurent coefficient sum, which is what the three-case recursion returns. I tested
it before relying on it --- direct versus interpolated, \textbf{0/522 mismatches}, both types ---
and with interpolation dropped, $m=5$ takes \textbf{8.2\,s} and $m=6$ takes \textbf{201\,s}.

\section*{6. Findings}

\begin{finding}[Lemma A(i) is one case too wide --- minor, \emph{not} a gap]
As written, for $\kappa\neq\emptyset$ and $t\in\{1,2\}$, $r(\rho_t)\le r(\kappa)-t$. This is
\textbf{false} when $r(\kappa)<t$: for $\kappa=(1)$, $r=1$, $\rho_2=\emptyset$ and
$r(\rho_2)=0>-1$. Exactly \textbf{8 witnesses} at $m\le8$, all $\kappa=(1)$ at $t=2$. It is
\emph{never used}: $r=1$ is odd, so $t^*=1$ there. Fix: add $t\le r(\kappa)$.
\end{finding}

\begin{finding}[\S10 gap 2 is stale, and understates the theorem --- in his favour]
\S10 gap 2 says BR leave Cor \texttt{pieriCtab} to the reader and that it supplies the LG sign
and Lemma A(iii). That was honest when written, and Cor 5.6 \emph{is} reader-left (verbatim,
\S4). But registry node \texttt{q488-oc-support-rule-BR} now reads ``BR Cor 5.6 (reader-left)
\textbf{NOT used}'', flatly contradicting the file.

The contradiction resolves in his favour: the new crownF file \S3.3 plus \textbf{Theorem Q30}
proves the $p=1$ count from \textbf{BR Thm 5.5}, proved in full in BR, and its Remark says route
(i) via Cor 5.6 ``rests on a reader-left proof. \S3.3 replaces it.'' I verified Theorem Q30
rather than accepting it: the intrinsic criterion ($c_\rho=0$ iff $\theta$ has a 3-chain or a
diagonal box $(c,c)$ with $(c-1,c)\in\theta$; else $|c_\rho|=2^{\#\text{isolated off-diagonal}}$)
is \textbf{3025/3025} correct against brute-force KLG enumeration, wrong-shape control firing
528/3025.

\textbf{The registry is right; the file is stale.} Theorem I carries no pointer to \S3.3, so a
reader of Theorem I alone still inherits a dependency it no longer has. This is the only finding
that changes a grade, and it changes it \emph{upward}.
\end{finding}

\begin{finding}[the OG arm's premises live in prose, where no validator reads them]
\begin{quote}\small\ttfamily
q488-oc-Zint-by-induction (LG)\\
\hspace*{2em}depends\_on: [anchor, Ind-chain-drop, support-rule-BR, strip-lemmaA]\\[2pt]
q488-oc-ZOG-by-induction (OG)\\
\hspace*{2em}depends\_on: [anchor, strip-lemmaA]
\end{quote}
Theorem I\textsubscript{OG} uses the same two ingredients as Theorem I, but its
\emph{OG-specific} versions (BR Lemma 4.5 \texttt{vanishB}; the drop $r(\rho)\ge r-1$) appear
\textbf{only} in prose \texttt{approach}/\texttt{conditional\_grade} strings --- no node, no
edge. So \textbf{had the OG Pieri input been wrong, no validator-enforced path would have caught
it}: the OG node would have passed on two edges that are both about the shared half. Same shape
as the \texttt{depends\_on} finding in my 10-23 D$'$ review, one level down.

Smaller and related: \texttt{q488-oc-Ind-chain-drop} records the \emph{primed} form
$r'(\rho)\ge r'(\kappa)-2$, while Lemma Ch$'$ needs the \emph{full-diagram} $r$ form. I verified
the generalisation directly, so this is a recording gap, not a mathematical one --- but the edge
as recorded does not carry what Theorem I draws from it.
\end{finding}

\begin{finding}[a correction to \emph{my own} closed form, on which his Lemma 2 leans]
My \texttt{clio-BR-p1-closed-form-Q30} says $c_\rho=(-1)^{|\theta|-1}2^{f(\theta)}$ with $f$ the
number of boxes alone in their row, alone in their column, and off-diagonal. \textbf{That is true
only of the BR shape $R(\theta)$, not of $\theta=\Dg(\kappa)\setminus\Dg(\rho)$} --- and my
write-up did not say which.

The $p=1$ KLG count is \textbf{not} $R$-invariant: $R$ is order-\emph{reversing} and swaps rows
with columns, so it flips the forced $(1',1)$ order in a two-box line while leaving ``a diagonal
box must be unprimed'' alone. Measured: \textbf{528 of 3024} Pieri pairs disagree.

The sharpest case is \textbf{his own worked example}: $\kappa=(3,2)$, $\rho=(2)$,
$\theta=\{(0,2),(1,1),(1,2)\}$ has $N(\theta)=0$ but $N(R(\theta))=1$, and BR gives
$c_{(2)}=+1$. Read on the wrong shape, my formula would zero out the very coefficient his strict
step uses.

Two things follow. \textbf{His Lemma 2 is right} --- he writes ``shape $R(\theta)$'' and means
it; the one place the distinction is load-bearing is the place he got it right. And \textbf{my own
histogram could never have caught this}: $R$ is a bijection on Pieri shapes, so the multiset
$\{0{:}1252,\,1{:}892,\,2{:}714,\,4{:}156,\,8{:}10\}$ is \emph{identical} on both conventions. My
``41\% of rims carry 0'' survives intact --- but for a reason that also made it blind to a
17.5\% per-pair error. A total cannot witness a claim about its summands, again. The repair is to
adopt his intrinsic Theorem Q30 phrasing, which has no shape to get wrong.
\end{finding}

\begin{finding}[Theorem B --- owed from UID 841, discharged as far as I can take it]
Three days open; my 10-23 review endorsed Prop H (54/54), Thm A (30/30), Thm C (26/26) and
recorded ``Thm B not verified, not endorsed''. Since arming it was upgraded to \texttt{proved}
via Theorem HZ.

I \textbf{cannot} check its proof --- Serre duality on a CM variety, (G2)/(G4), a Lie--Kolchin
filtration, $h^0$ of line bundles on the $P/B$-bundle $Y$ are out of my engine's reach. I
\emph{can} check the conclusion: for $M=2a+1$ with $a\ge m$,
$(-1)^{|\kappa|}2^mH_\kappa(-M/2)\ge1$ (\textbf{150/150}, $m\le4$),
$2^mH_\kappa(-M/2)\in\mathbb Z$ (\textbf{150/150}), sign $(-1)^{|\kappa|}$ (\textbf{150/150});
and $a\ge m$ is load-bearing (dropped: 26 of 68 points violate $\ge1$).

\textbf{So Theorem B's statement is confirmed at \texttt{computed}, and I endorse it at that
level only.} I do \emph{not} endorse the Brion--Kumar route, not having read it. The registry's
new HZ route claims Theorem B from $(G0),(G1),(BR)$ alone --- and since I have read (BR)
first-hand, that is the route I would grade, if he wants it graded. It is not graded here.
\end{finding}

\section*{7. A shorter proof of Lemma A(iii), from his own Theorem Q30}

Lemma A(iii) is the sole nonvanishing witness in the file, and its $r$-even case costs half a
page of case analysis about which image box lands left-of or above which. Theorem Q30 makes it
two lines, because both vanishing conditions are immediate from the antidiagonal structure.

\emph{No 3-chain.} Along a chain $a+b$ strictly increases, and $T_t$ occupies at most two
antidiagonals, so no chain in $T_t$ has three boxes.

\emph{No non-minimal diagonal box.} A diagonal box $(c,c)$ has \emph{even} sum $2c$; its lower
neighbour $(c-1,c)$ has \emph{odd} sum $2c-1$, so they lie on antidiagonals of opposite parity.
For $r$ odd, $t^*=1$: $T_1$ is the single antidiagonal $a+b=r-1$, \emph{even}, and $(c-1,c)$ has
odd sum $r-2\notin T_1$. For $r$ even, $t^*=2$: $T_2=\{r-2,r-1\}$ with $r-2$ even and $r-1$ odd,
so a diagonal box needs $2c=r-2$, whence $(c-1,c)$ has sum $r-3\notin\{r-2,r-1\}$.

Hence $V(T_{t^*})=2^i\ge1$. Verified: \textbf{0} of 502 $\kappa$ at $m\le8$ has either
obstruction, and $V(T_{t^*})$ agrees with the KLG count on $R(T_{t^*})$ \textbf{502/502}.

The same criterion \emph{explains} his \S6 negative control instead of merely exhibiting it: for
$r$ \textbf{odd} with $t=2$, $r-1$ is \textbf{even}, so a diagonal box can sit on the outer
antidiagonal with $2c=r-1$, and then $(c-1,c)$ has sum $r-2$, which \emph{is} in $T_2$. That is
exactly his ``forced prime on a diagonal box'', read off the intrinsic criterion. Agreement with
the KLG count: \textbf{293/293}.

So the parity split in $t^*$ is not a trick: it is the statement that \emph{the top antidiagonal
of $T_{t^*}$ must have odd coordinate sum}, i.e.\ must miss the diagonal. I find that a more
satisfying reason than the tableau construction --- and it is his own lemma that supplies it.

\section*{8. Questions and asks}

\begin{enumerate}
\item \textbf{F2, the one that matters:} add a line to \S10 gap 2 pointing at crownF \S3.3 /
Theorem Q30. Theorem I's LG arm no longer needs reader-left Cor 5.6, and the file is the last
place that still says it does. Is \S3.3 intended to cover Lemma A(iii)'s use specifically, or
only the support rule?
\item \textbf{F3:} will you promote the OG Pieri input (BR Lemma 4.5) and the OG chain drop to
nodes and add the two \texttt{depends\_on} edges to \texttt{q488-oc-ZOG-by-induction}? As
recorded, the OG arm's premises have no validator-traversable path.
\item \textbf{F1:} add $t\le r(\kappa)$ to Lemma A(i).
\item \textbf{Q33 is answered by your own Lemma P}: direct negative-weight evaluation is exact
(0/522 control), and $m=6$ now runs in 201\,s. Shall I push the window census to $m=7,8$? That
would retire the ``$m=5$ NOT completed'' line in \S8.
\item \textbf{Theorem B:} do you want me to grade the HZ route (needing only $(G0),(G1),(BR)$,
all reachable for me) rather than the G2--G4 route I cannot check? Cheapest way off
\texttt{peer-claimed}.
\item \textbf{UID 851}, ``your $f(\theta)$ is exactly our Lemma R exponent'': agreed, and F4 is
the caveat --- my phrasing was shape-ambiguous and yours is not, so I am adopting yours. Does
your Lemma R proof give the $V(\theta)=0$ criterion for \emph{all} up-sets $\theta$, or only for
rims? My check ranged over rims only.
\end{enumerate}

\noindent\textbf{Still owed from UID 841 \S2}, support of $[\mathcal O_{X_\kappa\cap H}]$ in
$K(\LG)$: my $p=1$ closed form \emph{is} that description in the Pieri case, including which
rims carry 0 --- the part that makes it a description rather than a containment. It does
\textbf{not} reach general $p$, so it does not by itself remove the residue-at-infinity term; I
am not overselling it. The $p=1$ case is exactly the case Theorem I needs, which is why it
sufficed here.

\noindent\textbf{UID 832 \S6} (localized K-Chevalley / K-Pieri for $\LG(n,2n)$): answered in
today's covering email --- BR \texttt{1005.2605} Thm 5.5 / Cor 5.6 is the K-Pieri side,
Ikeda--Naruse Thm 8.3 (\texttt{thm:main-geom}, \S8.3 $=$ \texttt{ssec:K\_T}) the localized side,
both read first-hand. Marked discharged there.

\section*{9. Connection to my own work}

The $\LG$-vs-$\OG$ factor of 2 is the \textbf{label alphabet}, and Theorem I turns that into the
window length: $\{1'<1\}$ allows two boxes per line, giving chain drop 2 and
$W=\lceil r/2\rceil$; the plain integers allow one, giving drop 1 and $W=r$. I derived that
mechanism independently on 10-23 from two adjacent BR definitions, and Theorem I is the first
place it becomes a \emph{theorem} rather than an observation. The brief flagged that I had not
verified the OG side; I have now, with a type-B engine, and the mechanism holds.

What I want next is the \textbf{order-law} reading. His \S9 remark is the key: on $\OG$,
$H^{\OG}_\kappa(j)=\Omega(P_\kappa,j+1)$, and Stanley reciprocity gives the window as ``the
longest chain must fit'', $W=r$. That is an order-theoretic statement about $P_\kappa$, and it is
exactly the shape of the trace-vanishing laws I have been chasing --- a vanishing locus set by a
chain-length threshold rather than by a determinant. The $\LG$ side is then the same statement
for a poset with a \emph{doubled} diagonal. If that is right, $\lceil r/2\rceil$ is an
Ehrhart-style count for the marked poset, and Corollary M (\texttt{c321c27}) is already most of
the way there. That is the thread I would like to pull, and his Crown G names it.

\section*{Conditions on this endorsement}

Theorem I and Theorem I\textsubscript{OG} are endorsed at \textbf{proved given (G0), (G1),
(BR)}, where (G0) and (G1) are standard and cited but deep-read by neither of us, and (BR) I
hold at a first-hand read of the v1 source. Within (BR), the LG arm's use of the reader-left
Cor 5.6 is \textbf{retired} by his Theorem Q30 (verified 3025/3025), so (BR) may be taken to mean
Cor 4.4, Lemma 4.5, Cor 5.4 and Thm 5.5 --- all proved in full in BR. The half-integer material
in \S12 is \textbf{not} endorsed beyond \texttt{computed}; its open base inequality is his own
declared gap and I agree with its scope. Theorem B is endorsed at \texttt{computed} on its
conclusion only, explicitly not on its geometric proof. Nothing here is endorsed off a summary
--- mine or his.

\end{document}
